% =============================================================================
% SECTION 4: COMPOSITION INFLATION AND THE RESIDUE AS CAUSAL PROPAGATOR
% =============================================================================

\section{Composition Inflation, the Causal Residue, and the Origin of the
Second}
\label{sec:composition}

\subsection{Composition Inflation}
\label{subsec:comp_inflation}

Each partition event is not just a state transition --- it is a branching of the
space of possible next states. The partition count $M$ advances by 1, but the
number of distinguishable ways it can advance depends on the number of available
partition directions $d$.

\begin{theorem}[Composition Inflation]
\label{thm:comp_inflation}
For a system with $d$ distinguishable partition directions, the number of
distinguishable categorical trajectories of length $n$ is:
\begin{equation}
    T(n,d) = d(d+1)^{n-1}
    \label{eq:comp_inflation}
\end{equation}
\end{theorem}

\begin{proof}
At depth 1, there are $d$ distinguishable starting directions: $T(1,d) = d$.
At each subsequent depth, each existing trajectory can branch into $d+1$
directions (the $d$ original directions plus the ``stay'' option, which
corresponds to remaining in the current partition state). Therefore
$T(n,d) = d(d+1)^{n-1}$.
\end{proof}

For $d=3$ (three spatial partition directions): $T(n,3) = 3 \cdot 4^{n-1}$.
For Cs-133 at $n_P = 56$: $T(56,3) = 3 \cdot 4^{55} = 3.81 \times 10^{33}$,
exceeding the $2.02 \times 10^{33}$ Planck intervals per Cs-133 oscillation
period. The partition framework generates more distinguishable trajectories than
there are Planck intervals, confirming that the Cs-133 oscillation is a
legitimate classical process fully resolved by the partition structure at
depth 56.

\subsection{The Residue as Causal Propagator}
\label{subsec:causal_residue}

We now establish the central claim of this paper.

\begin{theorem}[Residue as Causal Propagator]
\label{thm:causal_residue}
The irreducible residue $\beta > 0$ of partition $k$ is the necessary and
sufficient condition for the existence of partition $k+1$. Without the residue,
there is no causal link between successive partitions, and the universe has no
mechanism to progress from one state to the next.
\end{theorem}

\begin{proof}
Partition $k$ establishes the distinction $\mathcal{I} \neq \overline{\mathcal{I}}$
and generates residue $\beta_k > 0$. This residue is deposited at the boundary
between $\mathcal{I}$ and $\overline{\mathcal{I}}$. The residue constitutes a
physical difference between the state of the system immediately before partition
$k$ and immediately after. Partition $k+1$ occurs at the boundary, driven by the
presence of $\beta_k$: the residue creates a non-equilibrium condition at the
boundary, which is the driving condition for the next partition event. If
$\beta_k = 0$, no physical difference would exist between the pre- and
post-partition states, no non-equilibrium condition would exist, and partition
$k+1$ would have no cause. The residue is therefore the causal link. It is
necessary (no residue implies no next partition) and sufficient (any non-zero
residue creates the driving condition for the next partition).
\end{proof}

\begin{corollary}[The Three Forms of the Residue]
\label{cor:three_forms}
The residue $\beta_k$ manifests in three forms depending on which aspect of
the next partition is being described:
\begin{enumerate}
    \item \textbf{As time}: $\beta_k$ determines the interval $\tau_{p,k+1}$
          of the next partition event. Time is accumulated residue of partitioning.

    \item \textbf{As entropy}: When many residues $\{\beta_k\}$ are present and
          their specific causal links to future partitions are not individually
          tracked, the aggregate is entropy. Shannon entropy
          $H = -\sum_k p_k \log p_k$ is the information content of the
          residue distribution when individual causal links are unresolved.

    \item \textbf{As information}: When the residue $\beta_k$ from partition $k$
          is the specific input to partition $k+1$ of a different item, the
          residue is \emph{information} --- the structural imprint of one
          partition on another's boundary conditions.
\end{enumerate}
Time, entropy, and information are not three different things. They are the
same residue described from three different perspectives on the next partition.
\end{corollary}

\subsection{The Angular Reformulation of Fundamental Constants}
\label{subsec:angular}

The composition inflation framework reveals the partition-theoretic origin of
fundamental constants. We state the key reductions:

\begin{align}
    c &= 2\pi \;\mathrm{rad/tick}
    \label{eq:c_angular} \\
    \hbar &= \frac{E_{\mathrm{tick}}}{2\pi}
    \label{eq:hbar_angular} \\
    k_B &= \frac{E_{\mathrm{tick}}}{\ln(d+1)}
    \label{eq:kb_angular}
\end{align}

where $E_{\mathrm{tick}}$ is the energy per partition tick (the minimum quantum
of action in the partition framework) and $d+1 = 4$ for three spatial dimensions.

\begin{theorem}[G is Irreducible]
\label{thm:G_irreducible}
The gravitational constant $G$ cannot be expressed in terms of $E_{\mathrm{tick}}$,
$d$, $\tau_P$, or any combination of partition counting quantities.
\end{theorem}

\begin{proof}
Every reducible constant ($c$, $\hbar$, $k_B$) describes the rate or quantum of
a partition process: $c$ is a propagation rate, $\hbar$ is an energy per radian,
$k_B$ is an energy per log-state. All are expressible as $E_{\mathrm{tick}}$
times a dimensionless counting factor. $G$ appears in the gravitational
acceleration $a = GM/r^2$, where $M$ is partition inertia (mass) and $r$ is
spatial distance (partition depth times $\lambda_P$). Substituting:
\begin{equation}
    a = G \cdot \frac{\mu/\alpha}{\lambda_P^2 \cdot D^2}
\end{equation}
For $a$ to be expressed in terms of partition counting quantities, $G$ would
need to reduce to $[E_{\mathrm{tick}}] \cdot [\mathrm{dimensionless}]$ divided
by the units of $\mu/(\alpha \lambda_P^2)$. But $\lambda_P$ is itself defined in
terms of $G$: $\lambda_P = \sqrt{\hbar G/c^3}$. Substituting this definition
into the expression for $a$ makes $G$ appear on both sides of the equation in a
form that does not cancel. No rearrangement eliminates $G$. It is therefore
irreducible.
\end{proof}

The irreducibility of $G$ reflects the partition-theoretic insight of
Remark~\ref{rem:G_irreducible}: $G$ is the contact geometry constant, not a
process constant. Contact --- the ground state of partition depth minimisation
--- is prior to the processes that approach it, and therefore $G$ is prior to
the constants that describe those processes.

\subsection{The Partition Second}
\label{subsec:partition_second}

\begin{definition}[Partition Second]
\label{def:partition_second}
The \emph{partition second} (Ps) is the elapsed duration expressed as a ratio
of partition counts:
\begin{equation}
    \Delta t \;[\mathrm{Ps}] = \frac{\Delta M_{\mathcal{A}}}{\dot{M}_{\mathcal{A}}
    \cdot \dot{M}_{\mathrm{Cs}}} = \frac{\Delta t \;[\mathrm{s}]}{1\;\mathrm{s}}
    \label{eq:partition_second}
\end{equation}
where $\mathcal{A}$ is any analyser and $\dot{M}_{\mathrm{Cs}} = 9{,}192{,}631{,}770$
Hz is the Cs-133 hyperfine rate.
\end{definition}

\begin{theorem}[Universal Conversion Constant]
\label{thm:K}
For an Orbitrap with geometry constant $\kappa$, the number of Cs-133 hyperfine
cycles per Orbitrap axial cycle at mass-to-charge ratio $m/z$ is:
\begin{equation}
    \frac{\dot{M}_{\mathrm{Cs}}}{\dot{M}_{\mathrm{Orb}}(m/z)} = K \cdot
    \sqrt{m/z}
    \label{eq:K_conversion}
\end{equation}
where:
\begin{equation}
    K = \frac{\dot{M}_{\mathrm{Cs}} \cdot 2\pi}{\sqrt{e\kappa/u}}
    = 588.016 \;\mathrm{Cs\;cycles\;Orbitrap\;cycle}^{-1}\;u^{-1/2}
    \label{eq:K_value}
\end{equation}
is a universal constant derivable from CODATA values alone.
\end{theorem}

\begin{proof}
From Eq.~\eqref{eq:mdot_orb}: $\dot{M}_{\mathrm{Orb}} = \sqrt{q\kappa/m}/(2\pi)$.
Therefore:
\begin{equation}
    \frac{\dot{M}_{\mathrm{Cs}}}{\dot{M}_{\mathrm{Orb}}} =
    \frac{\dot{M}_{\mathrm{Cs}} \cdot 2\pi}{\sqrt{q\kappa/m}} =
    \frac{\dot{M}_{\mathrm{Cs}} \cdot 2\pi}{\sqrt{e\kappa/u}} \cdot \sqrt{m/z}
    = K\sqrt{m/z}
\end{equation}
where we have used $q = e \cdot z$ and $m = u \cdot (m/z)$ with $e$ the
elementary charge and $u$ the atomic mass unit.
\end{proof}

\begin{corollary}[Analyser-Independent Time]
\label{cor:analyser_independent}
For any analyser $\mathcal{A}$ with partition rate $\dot{M}_{\mathcal{A}}$:
\begin{equation}
    \Delta t \;[\mathrm{s}] = \frac{\Delta M_{\mathcal{A}}}{\dot{M}_{\mathcal{A}}}
    \label{eq:analyser_independent}
\end{equation}
holds independently of which analyser is used. Two analysers with different
$\dot{M}$ accumulate different $\Delta M$ over the same duration, but the ratio
$\Delta M_{\mathcal{A}}/\dot{M}_{\mathcal{A}}$ is identical. The partition
second is analyser-invariant.
\end{corollary}

\subsection{The Analyser-Switching Time Jump}
\label{subsec:time_jump}

\begin{definition}[Partition Rate Discontinuity]
\label{def:rate_discontinuity}
When an ion transitions from analyser $\mathcal{A}_1$ to analyser $\mathcal{A}_2$
at physical time $t_s$, the partition rate changes discontinuously:
\begin{equation}
    \dot{M}(t) = \begin{cases}
        \dot{M}_{\mathcal{A}_1} & t < t_s \\
        \dot{M}_{\mathcal{A}_2} & t > t_s
    \end{cases}
    \label{eq:rate_discontinuity}
\end{equation}
The partition count $M(t)$ is continuous at $t_s$ but its rate is not.
\end{definition}

\begin{definition}[Time Jump]
\label{def:time_jump}
The \emph{time jump} $\Delta M_{\mathrm{jump}}$ is the difference in partition
count accumulated over interval $[t_s, t_s + \tau]$ between the two analysers:
\begin{equation}
    \Delta M_{\mathrm{jump}}(\tau) = \left(\dot{M}_{\mathcal{A}_2} -
    \dot{M}_{\mathcal{A}_1}\right) \cdot \tau
    \label{eq:time_jump}
\end{equation}
The time jump is the measurable signature of the metric discontinuity at the
analyser boundary --- the residue of the partition between the two analyser
regimes.
\end{definition}

\begin{theorem}[Time Jump Scaling]
\label{thm:jump_scaling}
For the Q-Orbitrap DDA cycle (quadrupole MS1, Orbitrap HCD), the time jump
$\Delta M_{\mathrm{jump}}$ for an ion at $m/z$ over MS1 window $t_{\mathrm{MS1}}$
is:
\begin{equation}
    \Delta M_{\mathrm{jump}}(m/z) = \left(\frac{\sqrt{q\kappa/m}}{2\pi} -
    \frac{q_u\Omega}{4\pi}\right) \cdot t_{\mathrm{MS1}}
    \label{eq:jump_scaling}
\end{equation}
This quantity scales as $(m/z)^{-1/2}$ for the Orbitrap term and is constant
for the quadrupole term. Therefore $\Delta M_{\mathrm{jump}}$ is a
monotone-decreasing function of $m/z$, larger for light ions and smaller for
heavy ions.
\end{theorem}

\subsection{Partition Gradient Trajectories}
\label{subsec:gradient_traj}

\begin{definition}[Partition Gradient Trajectory]
\label{def:gradient_traj}
For an ion at $m/z$ observed at retention times $\{t_k\}$ along a
chromatographic elution, the \emph{partition gradient trajectory} is the sequence:
\begin{equation}
    \{M_{\mathcal{A}}(t_k)\} = \{\dot{M}_{\mathcal{A}}(m/z) \cdot t_k\}
    \label{eq:gradient_traj}
\end{equation}
for analyser $\mathcal{A}$. Different analysers produce different trajectories
for the same ion at the same retention times, with slopes $\dot{M}_{\mathcal{A}}(m/z)$.
\end{definition}

\begin{theorem}[Retention Time from Partition Counts]
\label{thm:rt_from_partition}
For two observations of the same ion at $m/z$ at retention times $t_i$ and
$t_j$, the ratio of retention times equals the ratio of accumulated partition
counts:
\begin{equation}
    \frac{t_i}{t_j} = \frac{M_{\mathcal{A}}(t_i)}{M_{\mathcal{A}}(t_j)}
    \label{eq:rt_ratio}
\end{equation}
This holds for any analyser $\mathcal{A}$, since the same $\dot{M}_{\mathcal{A}}$
cancels from numerator and denominator.
\end{theorem}

\begin{corollary}[DDA Cycle Count from Partition Accumulation]
\label{cor:dda_cycles}
The number of complete DDA cycles elapsed before retention time $t$ is:
\begin{equation}
    n_{\mathrm{cycles}}(t) = \frac{t}{t_{\mathrm{cycle}}} =
    \frac{M_{\mathrm{total}}(t)}{M_{\mathrm{cycle}}}
    \label{eq:dda_cycles}
\end{equation}
where $M_{\mathrm{total}}(t) = \bar{\dot{M}} \cdot t$ is the total partition
count at time $t$ using the time-averaged partition rate $\bar{\dot{M}} =
\dot{M}_{\mathrm{quad}} \cdot f_{\mathrm{quad}} + \dot{M}_{\mathrm{Orb}} \cdot
f_{\mathrm{Orb}}$, and $M_{\mathrm{cycle}} = \bar{\dot{M}} \cdot t_{\mathrm{cycle}}$.
The number of analyser-switching events is therefore directly readable from the
partition accumulation without reference to any clock.
\end{corollary}
